\documentclass[a4paper,11pt]{article}
\usepackage[T1]{fontenc}
\usepackage[utf8]{inputenc}
\usepackage[top=2.5cm,bottom=2.5cm,left=2.5cm,right=2.5cm]{geometry}
\usepackage{hyperref}
\usepackage{xcolor}
\definecolor{accent}{HTML}{1F4E78}
\hypersetup{colorlinks=true, urlcolor=accent}
\pagestyle{empty}
\setlength{\parindent}{0pt}
\setlength{\parskip}{10pt}

\begin{document}

Gentile team di Röchling Automotive,

mi candido per la posizione di Software Developer per la sede di Laives. Ho tre anni di esperienza continuativa come sviluppatore full-stack, con esperienza diretta di migrazione di piattaforme legacy, integrazione API e database relazionali, il lavoro al centro del progetto di digitalizzazione che descrivete.

Per FlowSpace ho gestito la migrazione automatizzata di un marketplace di spazi di lavoro da uno stack legacy a Next.js e Laravel, mappando la struttura dati esistente sulla nuova piattaforma ed eseguendo il cutover mantenendo il prodotto in produzione, portandolo da 2.000 a oltre 4.500 utenti. Ho inoltre sviluppato le integrazioni API con OfficeRnD e Nexudus per le prenotazioni automatizzate.

Per Urlink ho costruito la piattaforma e la sua logica con sviluppo assistito da AI, usando Next.js e Supabase, che si appoggia su PostgreSQL, lo stesso tipo di database relazionale richiesto per questo ruolo.

Non ho esperienza diretta con Django o Vue, ma lavoro regolarmente con framework backend e frontend analoghi (Laravel, Next.js/React) e con ambienti containerizzati Docker, e mi adatto rapidamente a nuovi stack quando il pattern di lavoro è lo stesso.

Vivo a Venezia, sono autorizzato a lavorare in Italia senza necessità di sponsorship e disponibile per la modalità ibrida a Laives. Resto a vostra disposizione per un colloquio.

Cordiali saluti,\\
Prabhavit Gupta

\end{document}
